%% ma: L12-B08-0164
%% lop: 12
%% chuong: 2
%% bai: 8
%% dang: 12.8.52
%% loai: TN4
%% muc: VD
%% dap_an: "C"
%% nguon: "CĐ2.4.pdf giáo viên gửi 27/10/2026 (M07, Lê Minh Khôi – Nguyễn Anh Khôi, Chương II Chủ đề 4 Các bài toán về tâm tỉ cự), Câu 14"
%% kien_thuc: "Bài 8 (biểu thức tọa độ của tích vô hướng), lớp 10 (giá trị nhỏ nhất của tam thức bậc hai)"
%% trang_thai: cho-duyet
%% ly_do: "Dạng mới 12.8.52 đề xuất, chờ giáo viên duyệt"
\begin{ex}
Trong không gian $Oxyz$, cho ba điểm $A(0;1;2)$, $B(0;1;0)$, $C(3;1;1)$ và $M$ là điểm di động thuộc mặt phẳng $(Oyz)$. Tính giá trị nhỏ nhất của biểu thức $T=\dfrac32\overrightarrow{MA}\cdot\overrightarrow{MB}+\dfrac12\overrightarrow{MB}\cdot\overrightarrow{MC}$.
\choice{$T_{\min}=\dfrac{49}{16}$}{$T_{\min}=-\dfrac{49}{16}$}{\True $T_{\min}=-\dfrac{49}{32}$}{$T_{\min}=\dfrac{49}{32}$}
\loigiai{$M(0;y;z)$: $T=2y^2-4y+2z^2-\dfrac72z+2=2(y-1)^2+2\left(z-\dfrac78\right)^2-\dfrac{49}{32}$.

$T_{\min}=-\dfrac{49}{32}$.}
\end{ex}
